\subsection{逻辑函数的化简}
\subsubsection{卡诺图}
以 4 变量为例:
\begin{table}[H]
    \caption{4 变量卡诺图与其对应的最小项}
    \centering
    \begin{tabular}{|c|c|c|c|c|} \hline
        \bfseries AB \textbackslash CD & \bfseries 00 & \bfseries 01 & \bfseries 11 & \bfseries 10 \\\hline
        \bfseries 00                   & $m_0$        & $m_1$        & $m_3$        & $m_2$        \\\hline
        \bfseries 01                   & $m_4$        & $m_5$        & $m_7$        & $m_6$        \\\hline
        \bfseries 11                   & $m_{12}$     & $m_{13}$     & $m_{15}$     & $m_{14}$     \\\hline
        \bfseries 10                   & $m_8$        & $m_9$        & $m_{11}$     & $m_{10}$     \\\hline
    \end{tabular}
\end{table}

\underline{性质}
\begin{itemize}
    \item 纵轴从真值表中最高位变量 (例如 $A$) 开始表示, 横轴以真值表中最低位变量 (例如 $D$) 结束表示;
    \item 相邻性:
          \begin{itemize}
              \item 几何相邻即为最小项逻辑相邻;
              \item 具有对称轴, 对称位置满足逻辑相邻;
              \item $n$ 变量卡诺图的最小项具有 $n$ 个逻辑相邻的最小项.
          \end{itemize}
    \item 横纵轴的二进制编码采用\textit{格雷码};
    \item 出现在函数式中的最小项应填 1, 其他默认为 0.
\end{itemize}

\subsubsection{一般的逻辑函数的化简}
\begin{enumerate}
    \item 作出卡诺图;
    \item 圈 1:
          \begin{itemize}
              \item 以矩形圈出仅含 1 的\textit{逻辑相邻}最小项;
              \item 面积为 2, 4, 8, ...
              \item 圈 1 时要求可以圈已经被圈过的 1, 但后一次圈要包含先前没被圈过的 1;
              \item \textit{矩形面积尽可能大, 且矩形数量尽可能少};
              \item 要圈完所有的 1;
          \end{itemize}
    \item 将这些被圈出的最小项合并, 去除其中发生变化的因子, 保留不变的因子;
    \item 整理, 得出最终结果 (与或式).
\end{enumerate}

可以根据观察, 判断取反函数是否更容易计算. 取反后将圈 0 而不是 1.

\subsubsection{含无关项的逻辑函数的化简}
\textbf{约束项}\quad 恒为 0 的最小项, 通常有如下表示:
\begin{equation*}
    m_{k_1}+m_{k_2}+\cdots+m_{k_p}=0.
\end{equation*}

\textbf{任意项}\quad 取值对函数输出无关的最小项, 通常有如下表示 ($d(\cdot)$):
\begin{equation*}
    Y=\sum m(\cdot)+d(\cdot).
\end{equation*}

约束项和任意项统称\textbf{无关项}.

在卡诺图中, 无关项被记作 ``X'', 在实际圈 1 或 0 时, 可以被任意看作 1 或者 0.

``X'' 的优先级高于 1.
